\section{Finite probability and operator tools}
\label{a:finite}

Throughout this appendix, \(V\) is finite, each \(S_x\) is a
nonempty finite set, and \(\Omega=\prod_{x\in V}S_x\).
An empty product has its one empty configuration.  For a strictly
positive probability \(\mu\) on \(\Omega\), use
\(\langle f,h\rangle_\mu=\sum_\omega\mu(\omega)
\overline{f(\omega)}h(\omega)\).  If \(U\subset V\), write
\(\mu_{-U}\) for the marginal on \(V\setminus U\), and define
\[
 (E_Uf)(a,b)=\sum_{t\in\prod_{x\in U}S_x}
       \frac{\mu(t,b)}{\mu_{-U}(b)}f(t,b),\qquad
 \operatorname{Var}_U f=E_U|f-E_Uf|^2.
\]
All denominators here are positive.  We use the single-site Dirichlet
form and Poincare-gap convention of Definition~\ref{d:framemodel}.
More generally, a Poincare gap at least \(c\) for an operator \(H\)
means \(\langle f,Hf\rangle_\mu\ge c\operatorname{Var}_\mu f\)
for every \(f\).  This convention includes one-point spaces.

\begin{lemma}\label{l:CE}
For a positive probability \(\mu\) on a finite product
\(\Omega=\prod_{x\in V}S_x\), with nonempty finite alphabets,
\(E_U\) is a self-adjoint orthogonal projection preserving constants
and \(\mu\). For real or complex \(f\),
\[
 \langle f,(I-E_U)f\rangle_\mu
   =\mu(\operatorname{Var}_U f)=\|f-E_Uf\|_\mu^2,\qquad
 \operatorname{Var}_\mu f=\mu(\operatorname{Var}_U f)
                          +\operatorname{Var}_\mu(E_Uf).
\]
A finite sum \(\sum_a(I-E_{B_a})\), allowing repeated blocks,
has quadratic form \(\sum_a\mu(\operatorname{Var}_{B_a}f)\).
\end{lemma}

\begin{proof}
Fix a complement configuration \(b\) and put
\(p_b(t)=\mu(t,b)/\mu_{-U}(b)\).
On this fibre, \(E_Uf\) is the constant \(p_b(f)\).  Consequently
\(E_U^2=E_U\), \(E_U1=1\), and grouping the finite sum by \(b\)
gives
\[
 \langle f,E_Uh\rangle_\mu
   =\sum_b\mu_{-U}(b)\overline{p_b(f)}p_b(h)
   =\langle E_Uf,h\rangle_\mu,
 \qquad \mu(E_Uf)=\mu(f).
\]
Thus \(E_U\) is the orthogonal projection onto the functions of
the complement of \(U\).  Orthogonality of \(f-E_Uf\) and \(E_Uf\)
proves
\[
 \langle f,(I-E_U)f\rangle_\mu
   =\|f-E_Uf\|_\mu^2
   =\sum_b\mu_{-U}(b)p_b(|f-p_b(f)|^2).
\]
This is the first asserted identity.  Since \(\mu(E_Uf)=\mu(f)\),
the decomposition
\(f-\mu(f)=(f-E_Uf)+(E_Uf-\mu(f))\) is also orthogonal.
Its squared norm is the total-variance identity.  Summing the first
identity over block labels gives the last assertion, including repeated
blocks.  The argument uses complex conjugates as displayed and therefore
applies to complex functions.  For \(U=\varnothing\) it gives \(E_U=I\);
for \(U=V\) it gives the projection onto constants, also when \(V\)
is empty.
\end{proof}

\begin{lemma}\label{l:finitecoupling}
For probabilities \(p,p'\) on a finite set,
\[
 d_{\rm TV}(p,p')=\tfrac12\sum_z|p(z)-p'(z)|
          =1-\sum_z\min\{p(z),p'(z)\}.
\]
There exists a coupling with disagreement probability equal to this
quantity; every coupling has at least that disagreement probability.
Marginalization cannot increase total variation.
For laws on configurations in a finite \(V\), a marginal coupling
on \(Q\subset V\) can be completed to a coupling of the full laws.
If its disagreement probability on \(Q\) is at most \(\varepsilon\),
the completion can have expected Hamming cost at most
\(|V\setminus Q|+|V|\varepsilon\).
If both laws dominate \(p_0\nu\) for a common probability \(\nu\)
and \(0\le p_0\le1\), their total variation is at most \(1-p_0\).
\end{lemma}

\begin{proof}
For each state \(z\) put \(a(z)=\min\{p(z),p'(z)\}\) and
\(t=\sum_z a(z)\).  The equality
\(|p(z)-p'(z)|=p(z)+p'(z)-2a(z)\), summed over \(z\), proves
the formula for total variation.  If \(t<1\), define a joint law by
\[
 \gamma(z,z')=a(z)\mathbf1_{\{z=z'\}}
       +\frac{(p(z)-a(z))(p'(z')-a(z'))}{1-t}.
\]
Its two marginals are \(p,p'\).  At each \(z\), at least one of
\(p(z)-a(z)\) and \(p'(z)-a(z)\) is zero, so the second term
has no diagonal mass.  Hence its disagreement probability is \(1-t\).
This formula includes \(t=0\); if \(t=1\), use only the diagonal
law \(\gamma(z,z)=p(z)=p'(z)\).  In any coupling, the mass at
\((z,z)\) is at most \(a(z)\), proving optimality.

For a map \(T\) to a finite set, summing over its fibres gives
\[
 \sum_b\left|\sum_{z:Tz=b}(p(z)-p'(z))\right|
       \le\sum_z|p(z)-p'(z)|.
\]
Taking \(T\) to be a coordinate restriction proves the marginalization
claim.

Let \(\gamma_Q\) be any coupling of the two \(Q\) marginals.
After sampling \((a,a')\) from \(\gamma_Q\), sample the remaining
coordinates independently with laws
\(p(\,\cdot\mid n_Q=a)\) and \(p'(\,\cdot\mid n'_Q=a')\).
A zero-probability conditioning value has zero probability of being
sampled; its conditional law may be chosen arbitrarily.  Summing out
the other component verifies both full marginals.  Pointwise in this
completion,
\[
 d_H(n,n')\le |V\setminus Q|
                      +|V|\mathbf1_{\{n_Q\ne n'_Q\}}.
\]
Taking expectations proves the requested bound.  Empty coordinate sets
and one-point configuration spaces are included in this construction.
Finally, domination by \(p_0\nu\) implies
\[
 \sum_z\min\{p(z),p'(z)\}\ge p_0\sum_z\nu(z)=p_0.
\]
The total-variation formula gives the last assertion, also for
\(p_0=0\) or \(p_0=1\).
\end{proof}

\begin{lemma}\label{l:surgery}
Let \(\nu\) couple two positive finite-volume specification laws.
Fix a set \(F\) in the volume. For every input pair, resample \(F\)
from any coupling whose first and second marginals are the two
respective conditional laws given their own complements, and keep
the complements fixed. The resulting law still couples the original
two laws. The chosen coupling may depend on the entire input pair.
A finite composition of such kernels has the same property.
\end{lemma}

\begin{proof}
Write the two original laws as \(\mu,\mu'\), and an input pair as
\((\omega,\omega')\).  For any test function \(f\) of the first
output alone, its conditional expectation after the operation is
\((E_F^\mu f)(\omega)\): by the stipulated marginal of the chosen
coupling this value does not depend on \(\omega'\), even when the
joint coupling does.  Integrating against \(\nu\) therefore gives
\[
 \sum_{\omega,\omega'}\nu(\omega,\omega')
                  (E_F^\mu f)(\omega)
       =\mu(E_F^\mu f)=\mu(f),
\]
where the last equality is Lemma~\ref{l:CE}.  The same calculation
with a function of the second output gives its law \(\mu'\).
Consistency identifies these conditional laws with the specification
laws on \(F\) with their respective complements fixed.  Applying the
calculation at each stage proves the assertion about finite compositions.
If \(F\) is empty, its unique resampling is the identity, so this case
is included.
\end{proof}

\begin{lemma}\label{l:density}
Let \(\pi\) be a positive product probability on a finite product,
and \(d\mu=\varrho\,d\pi\), with \(0<m\le\varrho\le M<\infty\).
Then for every \(f\),
\[
 \mathcal D_\mu(f)\ge m\mathcal D_\pi(f),\qquad
 \operatorname{Var}_\mu f\le M\operatorname{Var}_\pi f,\qquad
 \mathcal D_\pi(f)\ge\operatorname{Var}_\pi f .
\]
Thus \(\lambda(\mu)\ge m/M\).
In a conditional frame block, absorbing site weights and every
fixed boundary factor into \(\pi\) gives
\(\lambda(\mu)\ge r_b^{|E(Q)|}\).
\end{lemma}

\begin{proof}
For any positive probability \(\nu\), Lemma~\ref{l:CE} gives the
following variational formulas, with the infima taken over complex
functions or constants when \(f\) is complex:
\[
 \nu(\operatorname{Var}_x f)
   =\min_{h:\,h\text{ depends only on }V\setminus\{x\}}
                      \nu(|f-h|^2),\qquad
 \operatorname{Var}_\nu f=\min_{z\in\mathbb C}\nu(|f-z|^2).
\]
Indeed the two minimizers are \(E_x^\nu f\) and \(\nu(f)\),
respectively; subtracting either minimizer splits the squared norm
orthogonally.  For every admissible \(h\), the lower density bound
implies \(\mu(|f-h|^2)\ge m\pi(|f-h|^2)\).
Taking the first minimum, and then summing over sites, gives
\(\mathcal D_\mu(f)\ge m\mathcal D_\pi(f)\).
Using the constant \(z=\pi(f)\) in the second minimum gives
\[
 \operatorname{Var}_\mu f\le\mu(|f-\pi(f)|^2)
                         \le M\operatorname{Var}_\pi f.
\]

For completeness, the product-law inequality follows from finite
orthogonal projections.  Enumerate \(V=\{x_1,\ldots,x_v\}\) and let
\(R_j=E_{x_j}^\pi\).  These projections commute because \(\pi\) is
a product probability.  Put \(f_0=f\) and
\(f_j=R_1\cdots R_jf\).  Then \(f_v=\pi(f)\).
The differences \(f_{j-1}-f_j\) are mutually orthogonal: for \(j<k\)
the former is killed by \(R_j\), while the latter is fixed by
\(R_j\).  Also
\[
 \|f_{j-1}-f_j\|_\pi
   =\|R_1\cdots R_{j-1}(I-R_j)f\|_\pi
   \le\|(I-R_j)f\|_\pi.
\]
Each projection has norm at most one by Lemma~\ref{l:CE}.
Adding the squared bounds proves
\(\operatorname{Var}_\pi f\le\mathcal D_\pi(f)\).
For \(v=0\) both sides vanish.  Combining all three inequalities
yields
\(\mathcal D_\mu(f)\ge(m/M)\operatorname{Var}_\mu f\).

We spell out the frame application.  By
Definitions~\ref{d:framemodel} and~\ref{d:frame}, and
Lemma~\ref{l:threeframes}, each conditional block law has the form
\[
 \mu(n)=\frac1{Z_Q}
       \prod_{x\in Q}h_x(n_x)
       \prod_{e=\{x,z\}\in E(Q)}(1+\varepsilon_e n_x\cdot n_z),
 \qquad h_x(t)>0.
\]
Here \(h_x\) is the site weight times all fixed boundary factors
incident to \(x\); the signs are \(+1\) for the plane law and
\(-1\) for the antiferromagnetic law.  Define
\(\pi_x(t)=h_x(t)/\sum_{a\in S_x}h_x(a)\) and
\(\pi=\bigotimes_{x\in Q}\pi_x\).
Proper colouring ensures that both endpoints of every internal or
boundary bond use distinct frames.  Thus all boundary factors are
positive, and every internal factor is between \(b\) and \(2-b\).
With \(e_Q=|E(Q)|\), set
\[
 \psi(n)=\prod_{e=\{x,z\}\in E(Q)}
                     (1+\varepsilon_e n_x\cdot n_z).
\]
Then \(d\mu/d\pi=\psi/\pi(\psi)\), and
\[
 \frac{b^{e_Q}}{\pi(\psi)}
      \le\frac{d\mu}{d\pi}
      \le\frac{(2-b)^{e_Q}}{\pi(\psi)}.
\]
The ratio of these bounds is \(r_b^{e_Q}\), proving the frame claim.
For an empty block the empty product is one and the Poincare inequality
is automatic.  Only internal bonds have entered this density ratio.
\end{proof}

\begin{lemma}\label{l:comparison}
Let \(\mu>0\) be on a finite product, and let \((B_a)\) be a
nonempty finite labelled family of nonempty blocks.
Suppose the single-site heat bath of every conditional law on each
\(B_a\) has Poincare gap at least \(g>0\), uniformly in its exterior fibre,
and \(\sum_a(I-E_{B_a})\) has Poincare gap at least \(c>0\).
If \(t_{\max}=\max_x\#\{a:x\in B_a\}\), then
\[
 \operatorname{Var}_\mu f\le
             \frac{t_{\max}}{cg}\mathcal D_\mu(f).
\]
Repeated blocks are counted with multiplicity.
\end{lemma}

\begin{proof}
Fix a block label \(a\) and a complement configuration \(b\).
Apply its assumed Poincare inequality to
\(f_b(t)=f(t,b)\) in the conditional probability
\(\mu(\,\cdot\mid n_{V\setminus B_a}=b)\).
If \(x\in B_a\), further conditioning on \(B_a\setminus\{x\}\)
fixes precisely all coordinates other than \(x\).  The one-site
conditional variance in this fibre is therefore the restriction of
\(\operatorname{Var}_x f\) for the original law.  Consequently,
\[
 \operatorname{Var}_{B_a}f
      \le\frac1g\sum_{x\in B_a}E_{B_a}
                                      (\operatorname{Var}_x f).
\]
Integrate and sum over labels.  Lemma~\ref{l:CE} and the finiteness
of all sums give
\[
 \begin{split}
 c\operatorname{Var}_\mu f
 &\le\sum_a\mu(\operatorname{Var}_{B_a}f)\\
 &\le\frac1g\sum_x\#\{a:x\in B_a\}
                         \mu(\operatorname{Var}_x f)
 \le\frac{t_{\max}}g\mathcal D_\mu(f).
 \end{split}
\]
Division by \(c>0\) proves the assertion.  Every sum is over labels,
so repetitions contribute their full multiplicity.  Conditional laws
on one-point spaces cause no exception because their variances vanish.
\end{proof}

\begin{lemma}\label{l:contraction}
Let \(P=A^{-1}\sum_aE_{B_a}\), where \(A\ge1\) is the number
of labels. Suppose \(0<c\le A\) and that for every two input
configurations at Hamming distance one, their \(P\)-output laws
have a coupling of expected Hamming distance at most \(1-c/A\).
Then \(\sum_a(I-E_{B_a})=A(I-P)\) has Poincare gap at least \(c\).
This applies to site-dependent finite alphabets and also to a
one-point product space.
\end{lemma}

\begin{proof}
If \(\Omega\) has one point, every variance is zero and the conclusion
holds.  Otherwise, for a real function \(f\) define
\[
 \operatorname{Lip}(f)=
   \max_{d_H(\omega,\omega')=1}|f(\omega)-f(\omega')|.
\]
Any pair can be joined by changing its disagreeing coordinates one at
a time.  This path has length \(d_H(\omega,\omega')\), so
\[
 |f(\omega)-f(\omega')|
            \le\operatorname{Lip}(f)d_H(\omega,\omega').
\]
In particular a nonconstant function has positive Lipschitz seminorm,
even if some site alphabets are singletons.  For an input pair at
distance one, apply the assumed coupling to this inequality and take
expectations.  With \(r=1-c/A\in[0,1)\), it gives
\[
 |Pf(\omega)-Pf(\omega')|\le r\operatorname{Lip}(f),
 \qquad \operatorname{Lip}(Pf)\le r\operatorname{Lip}(f).
\]

Lemma~\ref{l:CE} shows that \(P\) is real, self-adjoint, preserves
constants and \(\mu\), and is a positive contraction on \(L^2(\mu)\).
Indeed it is an average of orthogonal projections.  The real subspace
\(\{f:\mu(f)=0\}\) is invariant.  On it choose a real orthonormal
eigenbasis: conjugation by
\(\operatorname{diag}(\sqrt{\mu(\omega)})\) makes the matrix real
symmetric, so finite-dimensional orthogonal diagonalization applies.
If \(Pf=\lambda f\) is a nonzero vector of that subspace, it is
nonconstant and hence
\[
 |\lambda|\operatorname{Lip}(f)
       =\operatorname{Lip}(Pf)\le r\operatorname{Lip}(f).
\]
It follows that \(A(1-\lambda)\ge A(1-r)=c\).
Expanding \(f-\mu(f)\) in this eigenbasis proves
\[
 \langle f,A(I-P)f\rangle_\mu\ge c\operatorname{Var}_\mu f.
\]
Constants are annihilated by \(A(I-P)\), and applying the real
inequality to the real and imaginary parts proves the complex version.
The argument also covers \(c=A\), when \(r=0\).
\end{proof}

